\documentclass[12pt,a4paper]{article}
\usepackage[margin=1in]{geometry}
\usepackage{graphicx}
\usepackage[table]{xcolor}
\usepackage{array}
\usepackage{float}

\definecolor{mypurple}{RGB}{200,20,200}
% Same color expressed in HTML hexadecimal:
\definecolor{mypurpleHTML}{HTML}{C814C8}
% xcolor mixture: 60 percent red and 40 percent yellow.
\colorlet{myblend}{red!60!yellow}

\begin{document}

\section{Lists}
Following are some technologies that can be used to format content:
\begin{itemize}
  \item \LaTeX
  \item HTML (+CSS)
  \item Markdown
\end{itemize}

You would most likely use \LaTeX\ in these scenarios:
\begin{enumerate}
  \item To write research papers.
  \item To embed mathematical equations within other software,
        such as Google Docs or HTML.
  \item To create CVs, project reports and presentations.
\end{enumerate}

There are three main categories of electronic markup:
\begin{description}
  \item[Presentational] Used by traditional word-processing systems.
  Binary codes embedded within text produce a WYSIWYG result,
  such as Microsoft Word.
  \item[Procedural] Markup instructions tell a program how to process
  the text. TeX is a classic example.
  \item[Descriptive] Markup labels parts of a document according to
  their meaning rather than their visual appearance. LaTeX and HTML5
  provide descriptive/semantic structure.
\end{description}

\section{Tables and Colors}
Code for Table~\ref{tab:languages} is in this subsection, but the float
may be placed on a separate float page. Change the table's placement
argument to test [t], [b], [h], or [!htbp]. The active [p] option below
asks LaTeX to place it on a float page.

\begin{table}[p]
\centering
\renewcommand{\arraystretch}{1.2}
\begin{tabular}{|>{\raggedright\arraybackslash}p{2.5cm}
                |>{\raggedright\arraybackslash}p{2cm}
                |>{\raggedright\arraybackslash}p{2cm}
                |>{\raggedright\arraybackslash}p{3.5cm}|}
\hline
\rowcolor{green}
Language & Requires compilation & Output & Use\\ \hline
C/C++ & Yes & Executable & To make apps\\ \hline
\LaTeX & Yes & PDF & To make documents\\ \hline
HTML/CSS/JS & No & Webpage & To make websites\\ \hline
\rowcolor{yellow}
2.5 cm & 2 cm & 2 cm & 3.5 cm\\ \hline
\end{tabular}
\caption{This table should be rendered on a special page.}
\label{tab:languages}
\end{table}

RGB represents red, green and blue channel values. The \texttt{RGB} model
in xcolor accepts integer channel values from 0 to 255.
The \texttt{rgb} model accepts normalized values from 0 to 1.
HTML hexadecimal notation uses \texttt{\#RRGGBB}, with two hexadecimal
digits per color channel. Thus RGB(200,20,200) corresponds to C814C8.

\[
\textcolor{mypurple}{T(n)}
=
\textcolor{myblend}{aT(n/b)+f(n)}
\]

\section{Figures}

\subsection{Pictures: Rotation and Cropping}
% Replace filenames with the exact extracted asset files.
% If images are not available, this will produce a compilation error.
% See README for image extraction instructions.
\begin{figure}[htbp]
  \centering
  % Set angle to the rotation that matches the intended reference.
  % angle=0 leaves the circular MNIT logo upright.
  \includegraphics[width=0.32\textwidth,angle=0]
    {assets/mnit_logo.png}
  \caption{MNIT logo}
  \label{fig:mnit-logo}
\end{figure}

\begin{figure}[htbp]
  \centering
  % trim order: left, bottom, right, top.
  % Adjust the top trim after visually comparing with the reference.
  \includegraphics[width=\textwidth,trim=0 0 0 65pt,clip]
    {assets/mnit_building.jpg}
  \caption{This is a floating image and can appear anywhere.}
  \label{fig:mnit-building}
\end{figure}

\subsection{Multiple Images}
This subsection includes two minipages with equal width.
\begin{figure}[htbp]
\centering
\begin{minipage}[t]{0.48\textwidth}
  \centering
  \includegraphics[width=\linewidth]{assets/mnit_logo.png}
  \caption{This is the first minipage}
  \label{fig:minipage-one}
\end{minipage}\hfill
\begin{minipage}[t]{0.48\textwidth}
  \centering
  \includegraphics[width=\linewidth,trim=0 0 0 35pt,clip]
    {assets/mnit_building.jpg}
  \caption{This is the second minipage}
  \label{fig:minipage-two}
\end{minipage}
\end{figure}

\end{document}
